\chapter{Magnetic Perturbations and Distribution-Function Engineering}
\label{ch:magnetic-control}

\section{Guiding-center motion and adiabatic invariants}

In a magnetic field that varies slowly over a gyro-orbit, a charged particle is described by guiding-center motion. The relativistic gyrofrequency and Larmor radius are
\begin{equation}
\Omega_{ce}=\frac{eB}{\gamma m_e},
\qquad
\rho_e=\frac{p_\perp}{eB}.
\label{eq:gyro}
\end{equation}
The first adiabatic invariant is
\begin{equation}
\mu=\frac{p_\perp^2}{2m_eB},
\label{eq:mu}
\end{equation}
with relativistic conventions varying by formulation. Conservation of $\mu$ causes parallel and perpendicular energy exchange as a particle moves through a field gradient. Curvature, grad-$B$, and electric drifts determine cross-field motion, while field-line topology controls parallel access to material boundaries \cite{hazeltine2003,freidberg2014,wesson2011}.

The crucial scale separation is between perturbation structure and orbit width. Low-energy magnetized electrons follow field lines closely. More energetic electrons have larger gyroradii, larger drift excursions, and often stronger resonance overlap. A perturbation can therefore be weak in magnetic amplitude yet selective in energy or pitch angle.

\section{Representation of a magnetic perturbation}

A convenient toroidal representation is
\begin{equation}
\delta\mathbf{B}(\psi,\theta,\phi,t)
=\sum_{m,n}\delta\mathbf{B}_{mn}(\psi)
\cos\!\left(m\theta-n\phi-\omega_{mn}t+\varphi_{mn}\right),
\label{eq:rmp}
\end{equation}
where $m$ and $n$ are poloidal and toroidal mode numbers. Resonance occurs near rational surfaces satisfying $q(\psi)=m/n$ for field-line perturbations, or more generally when the particle frequencies satisfy
\begin{equation}
 l\omega_b + m\omega_\theta-n\omega_\phi-\omega_{mn}=0,
\label{eq:particle-resonance}
\end{equation}
with bounce, poloidal, and toroidal transit frequencies. In mirrors or linear systems, the analogous condition couples the perturbation to bounce and gyro motion.

The perturbation objective in this thesis is not necessarily global stochasticity. A broad stochastic layer can degrade bulk confinement and overload plasma-facing components. The desired state is selective phase-space transport: a loss or diffusion rate that is weak for the thermal bulk and rises for the radiatively expensive part of the distribution.

\section{Open-field-line escape}

If a perturbation connects confined orbits to open field lines, the simplest loss time is
\begin{equation}
\tau_{\parallel}\sim\frac{L_{\parallel}}{|v_{\parallel}|},
\label{eq:parallel-loss}
\end{equation}
where $L_{\parallel}$ is connection length. Energy selectivity enters because high-energy electrons traverse the connection rapidly and may have orbit widths sufficient to sample the open region. A phenomenological loss operator is
\begin{equation}
L_{\mathrm{open}}[f]
=\nu_{\mathrm{open}}(E,\xi,\mathbf{x})f,
\qquad \xi=\frac{p_\parallel}{p}.
\label{eq:open-operator}
\end{equation}
One useful parameterization is
\begin{equation}
\nu_{\mathrm{open}}(E,\xi)=
\nu_0\left[1+\left(\frac{E}{E_*}\right)^a\right]
\mathcal{A}(\xi),
\label{eq:energy-loss-rate}
\end{equation}
where $E_*$ is the onset energy, $a$ controls selectivity, and $\mathcal{A}(\xi)$ selects the affected pitch-angle region. Equation~\eqref{eq:energy-loss-rate} is a model to be calibrated from orbit following, not a universal law.

A practical design must intercept escaping electrons safely. If an electron loses confinement but strikes a high-$Z$ wall, the wall itself can become a bremsstrahlung converter. Source suppression in the plasma can then be replaced by localized target bremsstrahlung. Low-$Z$ sacrificial collectors, electrostatic deceleration, energy recovery, or magnetic expansion are therefore part of the physical closure.

\section{Mirror loss cones and ambipolar potentials}

In a magnetic mirror with ratio $R_m=B_{\max}/B_{\min}$, a particle is magnetically trapped if
\begin{equation}
\sin^2\alpha_0>\frac{1}{R_m},
\label{eq:losscone}
\end{equation}
where $\alpha_0$ is pitch angle at the field minimum. Electrostatic potential modifies this boundary. For electrons, a confining ambipolar potential can be several electron temperatures, so escaping electrons may be relativistic even when the bulk is only mildly relativistic. Ochs, Munirov, and Fisch derived relativistic confinement corrections and showed that the classical scaling must be modified when $T_e$ and $|e\phi|$ approach $m_ec^2$ \cite{ochs2023mirror,pastukhov1974}.

A combined mirror-electrostatic loss condition can be expressed schematically as
\begin{equation}
\frac{p_\parallel^2}{2\gamma m_e}
>\mu\Delta B+|e\Delta\phi|,
\label{eq:mirror-potential}
\end{equation}
with exact relativistic treatment based on conserved total energy and magnetic moment. Magnetic perturbations can scatter selected particles into the loss cone. If pitch-angle diffusion coefficient $D_{\xi\xi}$ increases with energy or resonance overlap, the effective confinement time falls above a controllable threshold.

\section{Stochastic magnetic transport}

Overlapping magnetic islands destroy nested flux surfaces and permit radial field-line diffusion. Rechester--Rosenbluth transport gives the classic scaling
\begin{equation}
D_{RR}\sim v_\parallel L_c\left(\frac{\delta B_\perp}{B_0}\right)^2,
\label{eq:RR}
\end{equation}
where $L_c$ is a correlation length \cite{rechester1978}. The corresponding radial escape time across width $a$ is roughly
\begin{equation}
\tau_{RR}\sim\frac{a^2}{D_{RR}}.
\end{equation}
Because $v_\parallel$ saturates near $c$, velocity alone provides limited energy selectivity at highly relativistic energy. Selectivity can still arise from orbit width, finite-Larmor-radius averaging, resonant phase-space structure, and the radial distribution of the energetic population.

For a desired cutoff, one may compare the loss time with collisional replenishment:
\begin{equation}
\tau_{\mathrm{loss}}(E_c)\sim\tau_{\mathrm{slow}}(E_c).
\label{eq:cutoff-balance}
\end{equation}
Above $E_c$, loss dominates and depletes the tail; below $E_c$, collisions maintain the bulk. This balance provides a physically interpretable definition of the effective cutoff.

\section{Resonant pitch-angle and momentum diffusion}

A time-dependent perturbation can exchange energy and momentum with electrons. Quasilinear theory gives a diffusion equation
\begin{equation}
\frac{\partial f}{\partial t}
=\frac{\partial}{\partial I_i}
\left(D_{ij}\frac{\partial f}{\partial I_j}\right),
\label{eq:ql}
\end{equation}
where $I_i$ are action variables. The diffusion tensor is concentrated near resonances such as Equation~\eqref{eq:particle-resonance}. A perturbation can therefore be designed to move electrons toward a loss cone, toward lower energy, or toward a region in which collisions rapidly thermalize them.

Pure magnetic fields do no work in the particle rest-frame description because $\mathbf{v}\cdot(\mathbf{v}\times\mathbf{B})=0$. Energy reduction by a ``magnetic perturbation'' is thus indirect. It arises from induced electric fields, wave-particle interactions, altered access to electrostatic potentials, collisional relaxation after pitch-angle redistribution, or physical escape. This distinction is important for energy conservation and should be stated explicitly in any experimental claim.

\section{Magnetic reconnection}

Reconnection changes field-line connectivity and can both accelerate and exhaust energetic particles. Contracting magnetic islands produce first-order Fermi acceleration, while parallel electric fields accelerate electrons near diffusion regions \cite{drake2006}. Reconnection is therefore not automatically a suppression mechanism. A controlled reconnection layer might expel superthermal electrons, but uncontrolled reconnection can create the very tail that drives bremsstrahlung.

The net effect is determined by a competition between acceleration $A_E(E)$ and escape $\nu_{\mathrm{esc}}(E)$ in an energy-space equation,
\begin{equation}
\frac{\partial f_E}{\partial t}
=-\frac{\partial}{\partial E}\left[A_E f_E-D_{EE}\frac{\partial f_E}{\partial E}\right]
-\nu_{\mathrm{esc}}f_E+C_E[f_E]+Q_E.
\label{eq:energy-fp}
\end{equation}
A suppressive configuration requires the combined operator to reduce the stationary high-energy population, not merely to create rapid turnover.

\section{Shock acceleration, double layers, and adiabatic losses}

Although not all are externally imposed magnetic perturbations, several natural mechanisms illustrate how cutoffs arise.

\subsection{Shock acceleration limits}
Diffusive shock acceleration competes with synchrotron, inverse-Compton, collisional, and escape losses. A maximum energy follows from
\begin{equation}
\tau_{\mathrm{acc}}(E_{\max})
=\min\left(\tau_{\mathrm{esc}},\tau_{\mathrm{syn}},\tau_{\mathrm{IC}},\tau_{\mathrm{coll}}\right).
\end{equation}
A perturbation that shortens escape or increases loss can lower $E_{\max}$, but in a fusion plasma it may also reduce useful confinement.

\subsection{Double layers}
A double layer produces a localized potential drop. Electrons below or above a threshold may be reflected, accelerated, or transmitted depending on sign and direction. The energy scale is simply
\begin{equation}
E_{\mathrm{DL}}\sim e|\Delta\phi|.
\end{equation}
Magnetic geometry can stabilize, localize, or connect double layers to open field lines. A double layer can create a sharply structured distribution, though it may also produce beam-driven instabilities.

\subsection{Adiabatic expansion}
For an expanding relativistic particle population, adiabatic losses satisfy approximately
\begin{equation}
\left(\frac{\dd E}{\dd t}\right)_{\mathrm{ad}}
=-\frac{1}{3}(\nabla\cdot\mathbf{u})p v.
\end{equation}
Magnetic nozzle expansion can therefore reduce electron energy before collection. This may suppress target bremsstrahlung if high-energy electrons are intercepted only after expansion and electrostatic deceleration.

\section{Self-generated magnetic perturbations}

An anisotropic distribution can drive Weibel or filamentation instabilities, producing magnetic fields that scatter the particles and reduce anisotropy \cite{weibel1959}. Self-generated turbulence can thus act as a feedback mechanism on the radiating distribution. It is not necessarily controllable, and the resulting filaments may create localized hot-electron channels. Nonetheless, measurements of fluctuation spectra and anisotropy are important because an externally imposed perturbation may interact constructively or destructively with the natural instability spectrum.

\section{Coupled relativistic Fokker--Planck model}

A practical axisymmetric model in momentum magnitude $p$ and pitch-angle cosine $\xi$ is
\begin{align}
\frac{\partial f}{\partial t}
={}&-\frac{1}{p^2}\frac{\partial}{\partial p}
\left[p^2\left(A_p f-D_{pp}\frac{\partial f}{\partial p}
-D_{p\xi}\frac{\partial f}{\partial\xi}\right)\right] \\
&+\frac{\partial}{\partial\xi}
\left[D_{\xi\xi}\frac{\partial f}{\partial\xi}
+D_{\xi p}\frac{\partial f}{\partial p}
-A_\xi f\right]
+Q-\nu_L(p,\xi)f.
\label{eq:fp-final}
\end{align}
Here $A_p$ includes collisional drag and radiation reaction; $D_{pp}$ and $D_{\xi\xi}$ include collisional and quasilinear diffusion; and $\nu_L$ is the orbit-loss rate computed from the perturbed fields. Synchrotron reaction is anisotropic and must be retained when $B$ is large. Bremsstrahlung reaction can also reshape runaway-electron distributions \cite{embreus2016}.

A self-consistent loop is required:
\begin{equation}
\delta\mathbf{B}
\rightarrow \nu_L,D_{ij}
\rightarrow f
\rightarrow \mathbf{j},P_{\mathrm{rad}},p_{\parallel,\perp}
\rightarrow \delta\mathbf{B}_{\mathrm{plasma}}.
\end{equation}
The plasma response may amplify, screen, rotate, or phase-shift the applied perturbation. Vacuum-field calculations alone are insufficient near resonances.

\section{Optimization criteria}

A candidate perturbation should be judged by a multi-objective function,
\begin{equation}
\mathcal{F}
=w_1\frac{P_{\mathrm{rad}}}{P_{\mathrm{rad},0}}
+w_2\frac{P_{e,\mathrm{wall}}}{P_{0}}
+w_3\frac{\Delta P_{\mathrm{fus}}^-}{P_{\mathrm{fus},0}}
+w_4\frac{\Delta W_{\mathrm{MHD}}}{W_0}
+w_5\frac{D_{\mathrm{bio}}}{D_{\mathrm{bio},0}},
\label{eq:multiobjective}
\end{equation}
where $P_{e,\mathrm{wall}}$ is charged-particle wall load, $\Delta P_{\mathrm{fus}}^-$ is lost fusion performance, $\Delta W_{\mathrm{MHD}}$ penalizes instability or confinement degradation, and $D_{\mathrm{bio}}$ is receptor dose. Weight choices express design priorities and should be accompanied by Pareto fronts rather than a single opaque optimum.

\section{Failure modes}

The principal failure modes are:
\begin{enumerate}
\item conversion of escaping electrons into bremsstrahlung at a wall;
\item increased synchrotron emission from enhanced perpendicular momentum;
\item bulk confinement degradation that overwhelms the radiative benefit;
\item induced runaway acceleration by perturbation electric fields;
\item local hot spots at strike points;
\item impurity influx and increased $Z_{\mathrm{eff}}$;
\item hardening of the residual photon spectrum;
\item unstable feedback between pressure anisotropy and magnetic response; and
\item diagnostic misinterpretation caused by anisotropic emission.
\end{enumerate}

\section{Summary}

Magnetic perturbations can generate an effective energy cutoff only through transport, resonance, electrostatic coupling, induced electric fields, collisional relaxation, or escape. The proposed control must be evaluated with a relativistic kinetic model and a closed energy balance. The next chapter assumes a calculated photon source and develops the transport and dosimetry required to assess biological consequence.
